\section{Endpoint polarization, vanishing, and the adjoint}\label{sec:endpoints}

\subsection{The same-charge endpoint calculation}

The physical scalar variations may be written
$\delta q_m=2\bar\zeta_{im}\Psi_i$ and
$\delta\bar q_m=2\bar\Psi_i\zeta_{im}$, with the related physical
generator conjugation; see \cite[Appendix~C]{Baker}. Complexification
does not make the actions on $q$ and $\bar q$ independently adjustable.
For the plus-normal choice the annihilated scalar components are
$q_2$ and $\bar q_1$.

Here is a basis-independent polarization version of that statement.
Suppress the nonzero spacetime/V factor of the charge, let its H factor
be $a\in\C^2\setminus\{0\}$, and write
\begin{equation}
 q(v)=\sum_m v_mq_m,\qquad\bar q(w)=\sum_m w_m\bar q_m,
 \qquad E=\begin{pmatrix}0&1\\-1&0\end{pmatrix}.
\end{equation}
In the generator convention just specified, the invariant tensors give
\begin{equation}\label{eq:endpoint-tensors}
 Q(a)q(v)\ \propto\ (a^TEv)\Psi,\qquad
 Q(a)\bar q(w)\ \propto\ (a^Tw)\widetilde\Psi.
\end{equation}
This follows by rotating the representative conditions
$Q_{i2}q_2=Q_{i2}\bar q_1=0$, with $E$ the fundamental antisymmetric
tensor and the second contraction the fundamental/dual evaluation.
The suppressed fermion components are arbitrary, so an operator
annihilator requires the scalar coefficient to vanish.

\begin{proposition}[Zero same-charge scalar pairing]\label{prop:endpoint-zero}
For the constant fixed-H-factor charge above, nonzero ordinary
fundamental and antifundamental scalar endpoints separately killed by
that charge have zero mutual free contraction. They are not a physical
adjoint pair.
\end{proposition}
\begin{proof}
The two nullspaces in \eqref{eq:endpoint-tensors} are
\begin{equation}
 v=\lambda a,\qquad w=\mu Ea.
\end{equation}
The H tensor of the free propagator is $\delta_{mn}$, so its polarization
factor is $v^Tw=\lambda\mu a^TEa=0$.
Physical conjugation instead requires $w=v^*$ and gives
$v^Tv^*=\norm v^2>0$. If $v=\lambda a\ne0$, then
$a^Tv^*=\lambda^*\norm a^2\ne0$, so the physical conjugate is not
killed by that same charge.
\end{proof}

In the representative basis, $q_2$ pairs nontrivially with $\bar q_2$,
whereas its same-charge partner is $\bar q_1$. The physical adjoint
pair is a useful nonzero control, but its dressed expectation still
requires an independent positivity argument. At distinct endpoints
the two supersymmetry variations contain independent fermion fields,
so they cannot cancel each other as bare operator identities. In
this ansatz the bulk line variation already vanishes separately.
Junction matter or a fermionic superconnection would change the
operator and require a new calculation.

\subsection{The full bilinear has nonzero R charge}

The fields in \eqref{eq:plane-connection} and \eqref{eq:all-connection}
are neutral under the H rotation fixing $X_6$. Assign $q_1,q_2$ weights
$+1/2,-1/2$. Then $\bar q_1$ also has weight $-1/2$ and
\begin{equation}\label{eq:charged-bilinear}
 \mathcal O_C=\bar q_1(b)U_M(C)q_2(a)
\end{equation}
has charge $-1$. If the vacuum and the definition of the operator
preserve this symmetry, its expectation obeys
\begin{equation}
 \vev{\mathcal O_C}=e^{-\ii\alpha}\vev{\mathcal O_C}
 \quad\text{for every }\alpha,\qquad\vev{\mathcal O_C}=0.
\end{equation}
This selection rule extends beyond the free contraction. It does not
construct an ultraviolet-finite line operator. A compensating charged
insertion, a different endpoint sector or a symmetry-breaking vacuum
changes its assumptions. A fixed-endpoint deformation of the same
neutral connection carries no compensating charge.

\subsection{Why the original semicircle is a nonzero control}

The restriction is to a constant Poincare charge and the displayed
tangent map. On \eqref{eq:semicircle}, that map has scalar direction
$n_8=-\sin\theta$, $n_6=\cos\theta$: the scalar normal direction
changes sign along with the tangent. The common Poincare charge
therefore keeps its H weight fixed.

The original semicircle with constant scalar arclength coupling instead
uses a position-dependent conformal Killing spinor. Its local endpoint
conditions can select different H factors at the two ends. Baker's
same-polarization semicircle has a nonzero leading endpoint contraction,
proportional to $g^2/(4\pi|a-b|)$ in that paper's normalization
\cite[Section~3.2]{Baker}. It is consequently a successful control
outside Proposition~\ref{prop:endpoint-zero}, rather than a
counterexample to it. This distinction is also why the rigidity theorem
and tangent-map construction must not be conflated.

\subsection{The adjoint changes the scalar map}

For a tangent-coupled line and Hermitian matrix fields,
\begin{equation}\label{eq:dagger}
 U_M(C)^\dagger=U_{-M}(C^{-1}),\qquad
 U_M(C^{-1})=U_M(C)^{-1}.
\end{equation}
To prove the first identity, reverse the product order in the discretized
ordered exponential and conjugate each exponent. The gauge term changes
sign under conjugation; after reversing the contour it is again
$\ii A_\mu\dd x^\mu$, while the scalar one-form has acquired $-M$.
The second identity is ordinary inverse transport. Thus contour
reversal with unchanged tangent map gives the inverse, not the
Hermitian adjoint. A reflected norm must implement the former identity
with the appropriate spacetime pullback. This is the purpose of the
next section.
